\chapter{Аналитическая часть}

\section{Большие языковые модели}

Большие языковые модели (БЯМ) --- это класс нейросетевых моделей,
обученных на огромных объёмах текстовых данных для понимания и генерации человеческого языка.
Они основаны преимущественно на механизме трансформеров и способны предсказывать следующее слово в последовательности,
что позволяет им создавать связный и контекстно релевантный текст.
Обучение таких моделей требует значительных вычислительных ресурсов и данных, собранных из разных источников.
Благодаря масштабу параметров БЯМ демонстрируют способность к выполнению широкого спектра задач: от ответов на вопросы и
перевода до написания кода и сочинения текстов.
Однако, БЯМ не обладают истинным пониманием и подвержены галлюцинациям, при недостатке инфорормации.


\section{Qwen3-Coder}

Qwen-Coder~\cite{qwen} --- это специализированная крупная языковая модель для генерации и понимания программного кода, разработанная Alibaba Cloud в рамках серии моделей Qwen. Модель обучена на обширном корпусе исходного кода из открытых репозиториев и поддерживает множество языков программирования, включая Python, C++, Java, JavaScript и другие. Qwen-Coder демонстрирует высокую эффективность в задачах автоматического завершения кода, генерации функций по описанию, рефакторинга и объяснения программных конструкций. Особое внимание уделено точности синтаксиса, семантической согласованности и способности работать с многофайловыми проектами. В отличие от универсальных моделей, Qwen-Coder оптимизирована именно для программистских задач.

\section{deepseek}

DeepSeek~\cite{deepseek} --- это серия крупных языковых моделей, разработанных китайской компанией DeepSeek AI,
ориентированных как на общие, так и на специализированные задачи, включая программирование,
математику и анализ технической документации. DeepSeek-Coder — версия, специально обученная на наборе исходного кода из репозиториев GitHub
и других открытых источников.

\section{ChatGPT 5.1}

ChatGPT 5.1~\cite{chat-gpt} --- это современная крупная языковая модель, разработанная компанией
OpenAI. Данная версия характеризуется улучшенными возможностями по логическому рас-
суждению, расширенной глубиной контекста и повышенной устойчивостью к неоднозначным
запросам. ChatGPT 5.1 оптимизирована для работы в интерактивных системах, где требуется
высокая точность, стабильность и способность к длительным диалогам. Модель демонстрирует
высокие результаты как в задачах генерации кода, так и в задачах анализа структурированных
и неструктурированных данных.

\section{Регулярные выражения}

Регулярные выражения --- это формальный язык для описания шаблонов текста, предназначенный для поиска, извлечения и проверки строк.
В языке Python работа с регулярными выражениями осуществляется с помощью стандартной библиотеки \texttt{re}~\cite{python-re},
которая предоставляет функции для сопоставления текста с заданным шаблоном,
выделения подстрок и анализа структуры данных.
С помощью специальных метасимволов можно задавать фиксированные последовательности,
альтернативы, повторения (квантификаторы) и привязки к позициям в строке,
что делает регулярные выражения гибким инструментом обработки текстов.
В лабораторной работе они используются для поиска некириллических символах в кириллических токенах.

\section{Вывод}

В аналитической части были рассмотрены основные понятия, необходимые для решения
поставленной задачи. Дано определение большим языковым моделям и описаны используемые
в работе нейросетевые системы Qwen Coder, DeepSeek, ChatGPT 5.1. Приведено описание
регулярных выражений и их роли в анализе текстовых данных.

